% -----------------------------------------------------------------------------
% Related Work
% -----------------------------------------------------------------------------
\section{Related Work}
\label{sec:related_work}
Research on machine-learning-based insider threat detection includes
engineered-feature classifiers, temporal sequence models, staged security
pipelines, teacher--student procedures, and differentiable tree
architectures. Table~\ref{tab:related_work_positioning} positions
representative model families according to their training relationship,
explanation basis, and relevance to the present study. The table provides a
focused architectural synthesis rather than an exhaustive taxonomy of
insider-threat research.

\begin{table*}[!t]
\caption{Positioning of Representative Model Families}
\label{tab:related_work_positioning}
\centering
\footnotesize
\renewcommand{\arraystretch}{1.15}
\setlength{\tabcolsep}{4pt}
\begin{tabularx}{\textwidth}
{@{}>{\raggedright\arraybackslash}p{2.35cm}
>{\raggedright\arraybackslash}p{2.70cm}
Y
Y
Y@{}}
\hline
\textbf{Family} &
\textbf{Representative work} &
\textbf{Training relationship} &
\textbf{Explanation basis} &
\textbf{Relevance and boundary} \\
\hline
Engineered-feature and anomaly models &
\cite{Grinsztajn2022,AlMhiqani2022,Jaiswal2024,LeZincirHeywood2021} &
Trained on engineered, aggregated, or anomaly-based features. &
Feature importance or anomaly scores. &
Order and spacing must be encoded in the inputs. \\
Temporal sequence and attention models &
\cite{YuanWu2021,Song2024,Pal2023,JainWallace2019Attention,
SerranoSmith2019Attention,WiegreffePinter2019Attention} &
Encoder and prediction head share one gradient path. &
Attention weights and time-step interventions. &
Attention magnitude is not faithful input importance. \\
Staged security pipelines &
\cite{He2024,Gayathri2024,Ketepalli2022LSTMAE} &
Separate stages; an encoder often feeds a separately trained classifier. &
Per-component, rarely integrated across stages. &
No jointly optimised sequence--differentiable-tree pathway. \\
Teacher--student learning &
\cite{Hinton2015Distillation,Gou2021KDSurvey,Furlanello2018BornAgain,
Shafee2020Mimic,Hsu2023TeacherStudent,Umair2025KDIDS} &
A teacher guides the student via outputs, representations, or soft targets. &
Teacher--student prediction or attribution comparison. &
Aimed at transfer, privacy, or compression, not route preservation. \\
Differentiable tree ensembles &
\cite{Kontschieder2015,popov2020node,Marton2024GRANDE} &
Differentiable routing lets trees and neural layers co-train. &
Routes, leaves, selected dimensions, and tree outputs. &
Traceability alone does not establish faithfulness. \\
Related structured tabular models &
\cite{Prokhorenkova2018CatBoost,ArikPfister2021TabNet,martins2016sparsemax} &
Oblivious boosting (CatBoost), attentive selection (TabNet), sparsemax
transform. &
Structured selection or feature weighting. &
Not an end-to-end temporal ODST head. \\
\hline
\end{tabularx}
\end{table*}

Engineered-feature and anomaly models remain demanding references when user
activity has already been converted into fixed or aggregated behavioural
features \cite{AlMhiqani2022,Jaiswal2024,LeZincirHeywood2021}. As noted in
Section~\ref{sec:introduction}, Grinsztajn et al. found tree ensembles
highly competitive with deep neural models on typical, balanced tabular
data \cite{Grinsztajn2022}; their benchmark excluded stream-like and
time-series tasks, so it supports the inclusion of demanding tree
baselines here without predicting their performance on severely imbalanced
chronological insider-threat sequences. The limitation of engineered-feature
models in this setting is therefore not weak classification capacity, but
the need to represent chronology and behavioural spacing explicitly before
or during training.

\smallskip
Temporal sequence models instead learn from ordered observations, while
attention mechanisms aggregate evidence across behavioural intervals
\cite{YuanWu2021,Song2024,Pal2023}. Security research also includes staged
architectures in which separate components perform feature preparation,
augmentation, preliminary detection, temporal analysis, or final
classification. He et al. present a double-layer enterprise-threat
detection framework, while Gayathri et al. use behavioural representation,
SPCAGAN-based augmentation, and Bayesian-neural-network anomaly detection
as successive stages \cite{He2024,Gayathri2024}. These studies demonstrate
the value of modular security pipelines, but they do not establish the
specific jointly optimised sequence-encoder--differentiable-tree
relationship examined in the present work. A comparable staged pattern --
a temporal encoder trained first, with its compressed representation then
used to fit a separately trained, non-differentiable classifier -- appears
in adjacent security domains: Ketepalli and Bulla train an LSTM
autoencoder to compress network traffic features and subsequently fit a
Random Forest classifier on the encoder's frozen output for network
intrusion detection \cite{Ketepalli2022LSTMAE}. This confirms the pattern
exists, but its application to insider-threat sequences specifically
remains, to our knowledge, unestablished in the reviewed literature.

\smallskip
Deep neural decision forests introduced probabilistic routing within a
neural computation graph \cite{Kontschieder2015}. Neural Oblivious Decision
Ensembles (NODE) extended this principle through differentiable oblivious
trees that can be trained end-to-end by backpropagation
\cite{popov2020node}. In an oblivious tree, nodes at the same depth use a
common feature-selection rule and threshold. The ODST head used in this
study adapts this principle by applying sparsemax selection to the learned
temporal representation, sigmoid routing around depth-wise thresholds, and
aggregation of leaf responses across multiple trees
\cite{popov2020node,martins2016sparsemax}. This implementation should
therefore be interpreted as an adapted differentiable oblivious-tree head
rather than an exact reproduction of every NODE operation.

\smallskip
Other structured tabular architectures are relevant but technically
distinct. CatBoost uses oblivious trees within a conventional
gradient-boosting procedure rather than as a neural classifier head
\cite{Prokhorenkova2018CatBoost}. TabNet performs sequential attentive
feature selection without using a tree ensemble
\cite{ArikPfister2021TabNet}, while Gradient-Based Decision Tree Ensembles (GRANDE) provides a separate
gradient-based decision-tree ensemble design
\cite{Marton2024GRANDE}. These approaches motivate structured and sparse
decision mechanisms, but are not interchangeable instances of the same
sequence--tree architecture.

\smallskip
Teacher--student learning provides a separate basis for optimisation.
Knowledge distillation commonly uses information from a trained teacher to
guide a student alongside supervision from the true labels
\cite{Hinton2015Distillation,Gou2021KDSurvey}. The student is not required
to be smaller than the teacher. Born-Again Neural Networks demonstrated
that a student may retain the same architecture while learning from a
previously trained model \cite{Furlanello2018BornAgain}.

\smallskip
Teacher guidance has also been investigated in security detection. Shafee
et al. transfer intrusion-detection knowledge from a teacher trained on
private data to a shareable student
\cite{Shafee2020Mimic}. Hsu et al. use multiple binary teacher classifiers
to guide a multi-label student for anomalous user-event activity associated
with insider threats \cite{Hsu2023TeacherStudent}. Umair et al. use
knowledge distillation to produce a smaller intrusion-detection student and
compare teacher and student feature attributions using SHAP
\cite{Umair2025KDIDS}. These studies establish that teacher guidance is
relevant to both intrusion detection and insider-threat-related
user-activity analysis. Their principal purposes, however, concern
knowledge transfer, privacy, multi-label learning, compression, or
lightweight deployment.

\smallskip
The present study uses a frozen same-architecture teacher for a different
optimisation purpose. The teacher constrains both prediction logits and
ODST routing probabilities while the complete bidirectional long short-term
memory (Bi-LSTM)--attention--ODST student is updated using the true labels.
This teacher-anchored logit-and-route regularisation is intended to limit
prediction and routing drift during end-to-end optimisation rather than to
compress the student, generate pseudo-labels, or enable
privacy-preserving sharing. The teacher is removed after training and does
not form part of the inference architecture.

\smallskip
The expression \emph{teacher-anchored} has also been used for upper-layer
alignment in language-model embedding distillation
\cite{Dao2026TALAS}. In that formulation, the teacher's final embedding
acts as a stationary reference for selected student layers. It is therefore
distinct from the present procedure, which constrains prediction logits and
ODST routing probabilities within a same-architecture
sequence--ensemble.

\smallskip
Explanation availability and explanation faithfulness must also be
distinguished. Attention weights can show where a trained model allocated
temporal emphasis, but their magnitudes do not necessarily correspond to
faithful original-input importance
\cite{JainWallace2019Attention,SerranoSmith2019Attention}. Other work has
argued that whether attention can serve as an explanation depends on the
definition of explanation and the diagnostic procedure used
\cite{WiegreffePinter2019Attention}. Attention should therefore be treated
as model-specific temporal evidence whose relationship to prediction must
be tested, rather than accepted automatically as a complete explanation.
Related concerns apply beyond attention: post-hoc feature-attribution and
saliency methods can also produce rankings that fail basic sanity checks
against the model or the data they claim to explain
\cite{Adebayo2018SanityChecks}, which motivates evaluating ODST
tree-contribution rankings through explicit intervention rather than
treating structural traceability as sufficient evidence on its own.

\smallskip
Original-input occlusion addresses a different level of evidence by
measuring prediction sensitivity when a behavioural variable is replaced
with a defined reference value. Native ODST outputs provide selected latent
dimensions, route probabilities, reached leaves, and individual tree
scores. These quantities make the internal decision process traceable, but
their availability alone does not establish that a ranking derived from
them is faithful.

\smallskip
Faithfulness concerns whether the reported explanation remains meaningfully
connected to the underlying prediction function
\cite{Dasgupta2022Faithfulness}. In the present study, this relationship is
evaluated through matched interventions: components ranked as influential
are altered and compared with appropriate randomly selected controls. The
empirical ODST faithfulness claim is consequently limited to
reference-centred tree-contribution rankings tested through
top-versus-random tree-output ablation. Routes and reached leaves are
reported as traceable internal outputs rather than as independently
validated faithful explanation levels.